\section{Positive operators}

As you can guess, in the dictionary we have been building (like where self-adjoint operator corresponds to real number), the idea of a positive operator corresponds to a positve real number. And the ordering built in this section is the usual ordering on the real numbers.
1
2
3 These are all rules we know are true in the real numbers, but which are true in this general context
4 The nice thing about positive operators is that an operator is positive if and only if there is a T such that the operator has the form TT*
5 Note that a similar thing holds for self-adjoint <-> real spectrum 
6
9
10
11 What famous result from real analysis is 9.3-3 generalizing?
12 Also here the converse holds that every  positive oprator is of the form T²
13
14 In other words, if T is infinitesimal, then any smaller value is too! 
15

\begin{exercise}{1}
Let $S$ and $T$ be bounded self-adjoint linear operators on a complex Hilbert space.
If $S \leq T$ and $S \geq T$, show that $S = T$.
\end{exercise}
\begin{proof}
IF $S \leq T$, then $T - S \geq 0$, so that $\brackets{(T - S)x, x} \geq 0$, likewise, $\brackets{(S -T) x, x}  \geq 0$, thus $\brackets{(T - S) x, x} = 0$, by Lemma 3.8-2, we have that $T - S = 0$, as required.
\end{proof} 

\begin{exercise}{2}
Show that $T_1 \leq T_2$ if and only if $\brackets{T_1 x, x} \leq \brackets{T_2 x, x}$ defines a partial order relation on the set of all bounded self-adjoint linear operators on a complex Hilbert space $H$, and for any such operator $T$, $T_1 \leq T_2$ implies $T_1 + T \leq T_2 + T$ and $\alpha T_1 \leq \alpha T_2$ for a nonnegative $\alpha$.
\end{exercise}
\begin{proof}
(Reflexivity) Clearly $T \geq T$ for all $T$.

(Antisymmetry) This is the content of the previous exercise.

(Transitivity) If $T_1 \geq T_2$ and $T_2 \geq T_3$, we have $\brackets{T_1 x, x} \geq \brackets{T_2 x, x}$ and $\brackets{T_2 x, x} \geq \brackets{T_3 x, x}$ for all $x$, then $\brackets{T_1 x, x} \geq \brackets{T_3 x, x}$ and $T_1 \geq T_3$.

For the second part, suppose $T_2 \geq T_1$, so that $\brackets{T_2 x, x} \geq \brackets{T_1 x, x}$, then $\brackets{T_2 x, x} + \brackets{Tx, x} \geq \brackets{T_1 x, x} + \brackets{Tx, x}$, which is the same as $\brackets{(T_2 + T) x, x} \geq \brackets{(T_1 + T) x, x}$ so that $T_2 + T \geq T_1 + T$.
The second statement can be proven in a similar way using (positive) scalar multiplication instead of addition.
\end{proof} 

\begin{exercise}{3}
Let $A, B, T$ be bounded self-adjoint linear operators on a complex Hilbert space $H$.
If $T \geq 0$ and commutes with $A$ and $B$, show that $A \leq B$ implies $AT \leq BT$.
\end{exercise}
\begin{proof}
Since $B \geq A$, then $B - A \geq 0$, furthermore, $T$ commutes with $B - A$, since $T(B - A) = TB - TA = BT - AT = (B - A)T$.
By 9.3-1, the product of positive operators is positive, so that $T(B - A) \geq 0$;
that is, $TB - TA \geq 0$ which is the same as $TB \geq TA$.
\end{proof} 

\begin{exercise}{4}
If $T: H \to H$ is a bounded linear operator on a complex Hilbert space $H$, show that $T T^\ast$ and $T^\ast T$ are self-adjoint and positive.
Show that the spectra of $T T^\ast$ and $T^\ast T$ are real and cannot contain negative values.
What are the consequences of the second statement for a square matrix $A$?
\end{exercise}
\begin{proof}
We have $\brackets{T T^\ast x, y} = \brackets{T^\ast x, T^\ast y} = \brackets{x, TT^\ast y}$, so that $T T^\ast$ is self-adjoint.
Furthermore, $\brackets{T T^\ast x, x} = \brackets{T^\ast x, T^\ast x} = \norm{T^\ast x}^2 \geq 0$, so that $T T^\ast$ is positive.
The fact that $\sigma(T T^\ast)$ is real follows from 9.2-1, and the fact that $\sigma(T T^\ast)$ is nonnegative follows from the fact that $\inf_{\norm{x} = 1} \brackets{T T^\ast x, x} \geq 0$ and 9.2-1.

The implication for a square matrix $A$, is that $\bar{A}^T A = A A = A^2$ has nonnegative eigenvalues.

The proof for $T^\ast T$ follows using the same strategy mutatis mutandis.
\end{proof} 

\begin{exercise}{5}
Show that a bounded self-adjoint linear operator $T$ on a complex Hilbert space is positive if and only if its spectrum consists of nonnegative real values only.
What does this imply for a matrix?
\end{exercise}
\begin{proof}
We proved the necessary condition in the previous exercise.
We prove the sufficient condition by contrapositive.
Thus, suppose $T$ is not positive, so that there exists an $y \in H$ such that $\brackets{Ty, y} < 0$.
Then $\inf_{\norm{x} = 1} \brackets{Tx, x} < 0$ (this holds at least for $y$ normalised), and by 9.2-3, such infimum is a spectral value;
that is, not all the spectrum is nonnegative.
\end{proof} 

\begin{exercise}{6}
Let $T: H \to H$ and $W: H \to H$ be bounded linear operators on a complex Hilbert space $H$ and $S = W^\ast T W$.
Show that if $T$ is self-adjoint and positive, so is $S$.
\end{exercise}
\begin{proof}
We proved the fact that $S$ is self-adjoint in exercise 9.1.5.
To see $S$ is positive, we have $\brackets{Sx, x} = \brackets{TWx, Wx} = \brackets{Ty, y} \geq 0$ because $T$ is positive, so that $S$ is both self-adjoint and positive, completing the proof.
\end{proof} 

\begin{exercise}{9}
Show that if $T \geq 0$, then $(I + T)^{-1}$ exists.
\end{exercise}
\begin{proof}
We have $(I + T)x = x + Tx$ so that $Tx = -x$.
Moreover, $\brackets{Tx, x} = -\brackets{x, x} = -\norm{x}^2 \leq 0$, and $\brackets{Tx, x} \geq 0$, so that $-\norm{x}^2 = 0$, which implies $x = 0$.
Since $(I + T)0 = 0$, 2.6-10 tells us that $I + T$ has an inverse, as required.
\end{proof} 

\begin{exercise}{10}
Let $T$ be any bounded linear operator on a complex Hilbert space.
Show that the inverse of $I + T^\ast T$ exists.
\end{exercise}
\begin{proof}
This is a Corollary to exercises 9.3.4 and 9.3.9, since we proved there that $T^\ast T$ is positive
\end{proof} 

\begin{exercise}{11}
Show that an illustrative example for Theorem 9.3-3 is given by the sequence $(P_n)$, where $P_n$ is the projection of $\ell^2$ onto the subspace consisting of all sequences $x = (x_n) \in \ell^2$ such that $x_j = 0$ for all $j > n$.
\end{exercise}
\begin{proof}
We have that $K = I$ is bounded and self-adjoint and $\brackets{(I - P_n)x, x} = (\sum \absoluteValue{x_i}^2)^2 \geq 0$, so that $P_n \leq I$.
The limit of $(P_n)$ is $I$ so that, as expected, the conclusion of the Theorem holds.
\end{proof} 

\begin{exercise}{12}
If $T$ is a bounded self-adjoint linear operator on a complex Hilbert space $H$, show that $T^2$ is positive.
What does this imply for a matrix?
\end{exercise}
\begin{proof}
We have $\brackets{T^2 x, x} = \brackets{Tx, Tx} = \norm{Tx}^2 \geq 0$, so that $T^2$ is positive.
For a matrix, this implies that its square is a positive semi definite matrix.
\end{proof} 

\begin{exercise}{13}
If $T$ is a bounded self-adjoint linear operator on a complex Hilbert space $H$, show that the spectrum of $T^2$ cannot contain a negative value.
What theorem on matrices does this generalise?
\end{exercise}
\begin{proof}
If we combine exercise 5 and exercise 12, we get that $T^2$ has a nonnegative spectrum.
For matrices, this implies that $A^2 = A^T A$ has only nonnegative eigenvalues.
\end{proof} 

\begin{exercise}{14}
IF $T: H \to H$ and $S: H \to H$ are bounded linear operators and $T$ is compact and $S^\ast S \leq T^ast T$, show that $S$ is compact.
\end{exercise}
\begin{proof}
Let $(x_n)$ be a bounded sequence.
We want to prove $(Sx_n)$ contains a convergent subsequence.
We have $\brackets{(SS^\ast - TT^\ast) x, x} \leq 0$ so that $\brackets{Sx, Sx} \leq \brackets{Tx, Tx}$ which in turn implies $\norm{Sx}^2 \leq \norm{Tx}^2$.
In particular, this holds for all $x_n$, so that $\norm{Sx_n} \leq \norm{Tx_n}$
Since $T$ is compact, then $(Tx_n)$ contains a convergent subsequence, say $(Tx_{n_k})$, with $Tx_{n_k} \to Tx$.
We have that $\norm{Sx_{n_k} - Sx}$

FILL
\end{proof} 

\begin{exercise}{15}
Let $T: H \to H$ be a bounded linear operator on an infinite dimensional complex Hilbert space $H$.
If there is a $c > 0$ such that we have $\norm{Tx} \geq c \norm{x}$ for all $x \in H$, show that $T$ is not compact.
\end{exercise}
\begin{proof}
By Theorem 9.1-2, if there exists $c > 0$ such that $\norm{T_0 x} = \norm{Tx} \geq c \norm{x}$, then $0 \in \rho(T)$.
But we know from 8.4-4, that if $T$ is a compact operator defined on an infinite dimensional space, then 0 is in the spectrum, thus, it cannot be the case that $T$ is compact, as required.
\end{proof} 
